\documentclass[11pt]{article}
\usepackage[margin=0.82in]{geometry}
\usepackage{amsmath,amssymb,booktabs,longtable,enumitem,xurl,hyperref}
\hypersetup{colorlinks=true,urlcolor=blue}
\setlength{\parindent}{0pt}
\setlength{\parskip}{0.5em}
\newcommand{\R}{\mathbb R}
\newcommand{\E}{\mathcal E}
\newcommand{\B}{\mathcal B}
\newcommand{\support}{\operatorname{supp}}
\title{Why the Observer-Driven Controller Is Infeasible:\\A Compatible Prediction and Terminal Design}
\author{Collision Avoidance Research}
\date{October 3, 2026}
\begin{document}
\maketitle

\textbf{Finding.} The decisive mismatch is between the mathematical objects
used by the implementation: open-loop uncertainty propagation, a terminal
core proved for nominal exact-state feedback, and RTI candidates treated as
though affine feasibility implied a transferable nonlinear witness. These
objects do not establish the same closed-loop behavior. The target interface
also discards available information before constructing its prediction set.
These are formulation problems. A better potential-field seed or faster
conic solver cannot repair a negative-radius cone or a zero-gradient hard
collision row.

This report develops a correction to the existing single-controller,
single-CLF architecture. It does not implement that correction or claim a new
scenario success rate. The new executable artifact is an offline terminal
compatibility probe. The fifty-meter campaign and its failure snapshots are
the evidence base; the fourteen scenarios were not rerun for this design study.

\section{Causal ordering of the failures}
The completed campaign in
\path{report/PERCEPTION_50M_CAMPAIGN_20261003.tex} tested seven encounters at
8 and 15 m/s with the actual noisy NRMM estimator. All fourteen stopped at
the first control call, with solver primal infeasibility and no executed
hold. All targets were initially within 50 m. The same estimated ego state
and its published error box still caused failure after removing the target
in every case. Independently, setting only ego uncertainty to zero left all
fourteen target-uncertainty problems infeasible. These are offline
counterfactual solves, not successful closed-loop controls.

The generated terminal condition is
\begin{equation}\label{eq:empty}
 \|L_f(e_M+J_M\Delta U)\|_2\le r_f-R_M.
\end{equation}
The weighted uncertainty reserve $R_M$ exceeds $r_f$ by factors 439--3620.
Its right side is decision independent and negative. No choice of input,
trust radius, slack objective, or initialization can satisfy this cone.
The radii are weighted norm units, not meters. In physical terms, the pure
longitudinal-speed half-width of the nominal core is only
$3.22\times10^{-5}$ m/s at 8 m/s and $1.29\times10^{-4}$ m/s at 15 m/s.

The target prediction receives an initial course radius of $\pi$ even
though all snapshots contain 41 position-history samples and an available
direction interval. Its position enclosure reaches 277.55--305.67 m over
4.4--4.8 s. The interface tightens its published components using history,
then builds the prediction set from the earlier coarse components.
It also replaces the fitted initial velocity magnitude with the same
15 m/s prior in all cases. These are separate losses of useful information.

The exact-state comparison recovered in eleven cases and stopped without a
control in three. Two failures involve damping toward a primary point whose
nonlinear mismatch is already excessive. In the third, an overlapping
potential-field seed generates the hard row $0\Delta U\le-0.20$.
Thus fixing the observer interface alone will not resolve all known failures.
Conversely, those three RTI failures do not explain the common first-frame
observer failure: the empty cone already proves it independently.

\section{Three different errors must remain distinct}
Write $x$ for the true ego state, $\hat x$ for its estimate, and $c$ for an
MPC nominal center. The errors
\[
 \epsilon=x-\hat x,\qquad d=x-c
\]
have different dynamics. An observer ISS bound constrains $\epsilon$ about
the future estimate. It does not automatically contract the set of future
true states about today's open-loop prediction. Future measurements and
their centers are not yet known. An uncontrolled target's uncertain speed
can therefore cause a growing position set even while its continuously
updated observer is stable.

For a causal policy within the same MPC,
\begin{equation}\label{eq:policy}
 u_i=v_i+K_i(\hat x_i-c_i),
\end{equation}
one linearization at the dynamically consistent anchor gives
\begin{align}
 c_{i+1}&=f_h(\bar x_i,\bar u_i)+A_i(c_i-\bar x_i)+B_i(v_i-\bar u_i),\\
 d_{i+1}&=(A_i+B_iK_i)d_i-B_iK_i\epsilon_i+w_i+r_i.
 \label{eq:error}
\end{align}
Here $r_i$ encloses the nonlinear approximation remainder and $w_i$ the
declared plant uncertainty. Consequently a sound tube satisfies
\begin{equation}\label{eq:tube}
 \E_{i+1}\supseteq(A_i+B_iK_i)\E_i
 \oplus(-B_iK_i)\E_i^{\mathrm{obs}}\oplus\mathcal W_i\oplus\mathcal R_i.
\end{equation}
The minus sign follows from $\hat x-c=d-\epsilon$. The observer term cannot
be omitted simply because $A+BK$ is stable. The current implementation has
neither this feedback response nor this observer-input contribution in its
generator recursion; it uses $G_{i+1}=A_iG_i$.

Equation~\eqref{eq:policy} is a description of one controller's predicted
response, not a second executable backup. Its initial center and carried
state must be consistent with the input actually issued. Setting
$c_0=\hat x_k$ makes the first feedback correction zero, but does not alone
prove that resetting the center each frame preserves the previous witness.
There are two rigorous realizations: retain a center and its tube together,
or establish the recentering and shifted-candidate argument for a simplified
constraint-tightening MPC. Merely substituting $A+BK$ into existing code
while leaving the policy and terminal proof unchanged does neither.

The primary literature provides a closely related construction: K\"ohler,
M\"uller and Allg\"ower derive coupled observer/tracking error bounds and
information-dependent terminal ingredients. Their Section IV-C also gives a
simplified tightening formulation that applies the optimized first input
directly; its recursive proof requires additional terminal robustness and
incremental-stability assumptions \cite{Kohler}. This supports an efficient
design direction, not a claim that their theorem already covers the present
RTI bicycle implementation.

\section{The terminal core must tolerate estimation error}
The existing construction shrinks a nominal ellipsoid until a local RK4
contraction enclosure is established. That can be appropriate for the
nominal model. With nonzero observation error, shrinkage eventually opposes
invariance: feedback can push an exactly trimmed vehicle outside an
arbitrarily small core.

\subsection{New numerical counterexample using the actual saved cores}
\path{scripts/probeObserverTerminalCompatibility.m} loads the saved head-on
formulations at both speeds. It starts the true vehicle exactly at the
straight cruise trim, takes previous input equal to trim input, and applies
the existing terminal gain to a perturbed estimate. Only the lateral-speed
estimate is perturbed. The successor uses the actual nonlinear RK4 vehicle
map. A Schur-complement quotient minimizes over rigid endpoint pose, so a
failure cannot be repaired merely by translating or rotating the core.

\begin{center}
\begin{tabular}{rrrr}
\toprule
Speed & Nominal radius & $|\epsilon_{v_y}|$ (m/s) & Successor norm / radius\\
\midrule
8 & 0.00024414 & 0.00046834 & 0.9046\\
15 & 0.00097656 & 0.00046834 & 0.2021\\
8 & 0.00024414 & 0.00468343 & 9.0462\\
15 & 0.00097656 & 0.00468343 & 2.0202\\
8 & 0.00024414 & 0.46834303 & 898.4516\\
15 & 0.00097656 & 0.46834303 & 196.2560\\
\bottomrule
\end{tabular}
\end{center}
Both signs give the same displayed norms. The middle perturbation is only
one percent of the controller's admitted initial lateral-speed radius, yet
one hold leaves every rigidly placed copy of the nominal core. The exact
trim with zero estimation error stays inside. The smallest perturbations
also remain inside: this is a quantitative incompatibility, not a claim that
every nonzero error destroys invariance.

The nominal closed-loop spectral radius is approximately 0.98757 at both
speeds. Thus nominal local stability does not resolve the uncertainty
incompatibility. This probe extrapolates the nominal feedback to the
admitted error box; it neither executes that feedback in a campaign nor
claims that each box point is reachable by the actual observer. A tighter
correlated observer set could remove some box points, which is another
reason to improve the set representation.

\subsection{Construct the radius together with the noise floor}
Use a transverse, possibly input-augmented error measure $q=\|L_fe\|$.
If the actual estimate-feedback map admits the verified bound
\begin{equation}
 q^+\le a q+Cq^2+d_f,\qquad a<1,
\end{equation}
then a sublevel radius $r$ is invariant under this certificate if
\begin{equation}\label{eq:radius}
 Cr^2-(1-a)r+d_f\le0
\end{equation}
and the whole set satisfies the input and model-domain conditions. The term
$d_f$ includes estimation-error feedback, trim defect and admitted plant
disturbance. For $C=0$, this certificate requires $r\ge d_f/(1-a)$.
For $C>0$, the quadratic must have nonnegative discriminant and $r$ must
lie between its two roots. These are conditions for the chosen bound, not
a necessary test for every possible invariant-set representation.

The appropriate endpoint is therefore a family over nominal transverse
state, tube size, observer-error envelope and input memory. Pose/path phase
remains free. A terminal inequality can still have the form
$\|L_fe_M\|+R_M\le r$, but now $r$, $R_M$ and the appended policy belong
to the same robust construction. Enlarging $r$ by hand or discarding $R_M$
would not establish this. If \eqref{eq:radius} cannot be established inside
the admissible domain, improve the observer bound, feedback/metric or local
enclosure; do not report an empty construction as a physical impossibility.

An ego invariant core does not itself establish future target safety.
Matching-time target interaction must remain in the joint terminal family
or be closed by an actually justified departure condition. The experimental
50 m rule activates risk per time/position; a first exit alone does not
exclude later re-entry by a constant-curvature target.

\section{Use the target information already available}
The observer interface should maintain a joint current information set:
\begin{equation}
 \Theta_k=\Theta_k^{\mathrm{observer}}
       \cap\Theta_k^{\mathrm{history}}
       \cap\Theta^{\mathrm{declared\ motion}},
\end{equation}
with consistent timestamps and frames. Recenter the resulting set about the
published estimate without changing its physical contents. In particular,
a tighter inertial heading interval cannot be copied into a relative-course
field without accounting for shared ego-yaw error. The intersection must
reach \path{predictionErrorSet}, not only an unused component publication.

For constant tangential acceleration $A$ and sideslip $\beta$, let
$\kappa=\sin\beta/l_r$, $s=V_0\tau+A\tau^2/2$, and
$\gamma_0=\psi_0+\beta$. Before any stop, target position has the compact map
\begin{equation}
 p(\tau)=p_0+s\,\operatorname{sinc}(\kappa s/2)
 \begin{bmatrix}\cos(\gamma_0+\kappa s/2)\\
                 \sin(\gamma_0+\kappa s/2)\end{bmatrix},
 \qquad\operatorname{sinc}(z)=\frac{\sin z}{z}.
\end{equation}
If the model holds a braking target at rest after zero speed, use its stopped
arc length beyond that event. Each possible target keeps the same parameters
throughout a prediction. The next frame may use revised estimates, as the
user's experimental contract requires.

Propagate the joint parameter set through this map. Collision tightening
only needs support in a selected separating direction, for example
\[
 h_i(n)=\sup_{\theta\in\Theta_k}
 n^\top\big(p_i(\theta)-\bar p_i\big),
\]
with an appropriate joint orientation contribution. A validated enclosure
of that support can be computed before the convex solve, using sensitivities
and bounded remainders on a small parameter domain. Sampling alone is not a
certified supremum. Preserving direction and common parameter dependence
avoids replacing an informed motion estimate by a huge circular obstacle.
Uncertainty must still grow when the available information genuinely allows
diverging target paths; observer convergence cannot remove that fact.

Initialization should retain the fitted velocity magnitude with a justified
fit-error bound. When history identifies acceleration or curvature, use a
model-consistent fit and its residual/noise enclosure; when it does not,
retain uncertainty instead of inventing parameter accuracy. There is no
need to introduce $\dot\beta$ or an additional target motion mode.

\section{Remove local convexification dead ends}
\subsection{Overlapping seeds need an informative separating row}
For fixed unit $n$, let the two centered vehicle footprints be $S_E,S_T$.
The signed support gap
\begin{equation}\label{eq:gap}
 g_n=n^\top(p_T-p_E)-\support(S_E,n)-\support(S_T,-n)
\end{equation}
is negative at some overlaps but retains position derivatives $-n,n$.
The condition $g_n\ge d_0-\xi$ is a sufficient separation row when its
right side is nonnegative; a positive violation slack is not a safety
certificate. Choose $n$ consistently with the shifted trajectory or APF
avoidance side and freeze it within the convex problem. Enclose rotation
and dynamics linearization errors on the allowed correction domain.

This avoids linearizing the flat zero ordinary-distance value at an
overlapping seed, which produced $0\Delta U\le-0.20$. The row is informative
even there. It does not guarantee that every chosen normal, trust domain or
physical encounter is feasible, and APF initialization remains a search seed.

\subsection{Damping has to approach the claimed reference}
With $U(\alpha)=U_P+\alpha(U_S-U_P)$, the limit as $\alpha\to0$ is $U_P$.
If $U_P$ is an affine primary solution with nonzero nonlinear prediction
defect, shrinking $\alpha$ need not drive that defect to zero. The observed
17.756 and 32.608 mm primary defects explain the two failed searches.
Convex feasibility of a segment and nonlinear consistency of its base are
different properties.

For a genuinely nonlinear-consistent anchor, Taylor error about that anchor
can be bounded by a quantity quadratic in the total correction. A feasible
segment also requires its base to satisfy the same convex constraints.
If no such base exists, damping is not a feasibility argument. Construct an
informative feasibility-restoration subproblem or carry a properly transferable
policy; do not assume a saved affine solution is already a nonlinear witness.
Restoration variables describe the search and are not extra physical safety
permissions.

The more coherent one-linearization design incorporates a sound model-error
reserve into the same tube and directional collision constraints before
solving. This replaces an unrelated fixed global pose-error cutoff with
the error quantities the constraints actually require. It needs valid local
remainders, including tire-regime transitions, the complete correction
domain and any claimed intersample coverage. Increasing the 10 mm cutoff
alone provides none of these properties.

\section{Implementation order and what would count as resolution}
\begin{enumerate}[leftmargin=*]
\item Repair information transfer and initialization in
\path{nrmmControllerErrorBounds.m} and
\path{nrmmEstimatorControllerAdapter.m}: preserve the intersected history
set and fitted velocity, in the correct frame. Re-evaluate prediction support
with the original sensor data, without using scenario truth to tune estimates.
\item Jointly redesign the propagation in \path{solvePredictiveControl.m}
and the core in \path{terminalContinuation.m}. Select the same causal
estimate-feedback policy, derive observer and tracking/remainder bounds, and
construct a compatible invariant information family offline. Establish
target-free noisy cruise before trying to attribute failures to avoidance.
\item Repair signed collision support rows and the RTI reference/budget
contract. Preserve a nonlinear-consistent shifted initialization; regenerate
the existing APF seed only when needed. A changed target estimate is admitted
in experiments; inherited-budget theorems remain conditional on their
mathematical update assumptions, without adding an assumption-failure gate.
\item Keep the one given-path CLF. In the presence of persistent uncertainty,
the appropriate recovery inequality is
\[
 V(x^+)-V(x)\le-\eta\ell(x)+\delta+\sigma(b),
\]
where $\delta$ is the optimized CLF slack and $\sigma$ a derived uncertainty
supply. Zero slack gives disturbance-dependent recovery; vanishing supply
can recover asymptotic nominal behavior on the established recovery domain.
Persistent measurement noise does not generally permit an exact true-state
zero-error guarantee. The path objective remains unchanged.
\item Reproduce the three exact-state failed frames, then rerun all fourteen
actual-observer closed loops until recovery or collision/no solved control.
Report physical clearance, nominal recovery and actual call latency.
Information-bound premise violations alone remain outside the experimental
failure definition. Evaluate several sensor seeds before claiming robustness.
\end{enumerate}

Precomputed gains, set shapes and local bounds can leave the online problem
as one RTI convex correction with inherited safety budget when applicable,
and the existing PCBF initialization/restoration solve when necessary. The
secondary CLF problem remains present at zero PCBF slack. Neither a second
CLF nor an executable backup is required by this proposal. A one-shot local
problem cannot be promised feasible for arbitrary initializations; passing
all fourteen scenarios and a 0.1 s maximum latency remain measurements to
obtain after implementation.

\section{Evidence and limits}
The campaign source commit is
\path{aa85322df33f65da055475499fa62b15a9e3a377}. The new terminal probe is a
deterministic offline test of saved core/gain data, with no random draws.
Its twelve perturbed holds and two exact-trim controls are recorded in
\path{report/OBSERVER_CONTROL_COMPATIBILITY_20261003/terminal-probe.json}.
The finite mathematical checks in \path{verifyObserverRobustConnection.m}
were rerun: all seventeen passed. They exercise the existing derivation;
they do not validate its constants over all nonlinear vehicle trajectories.
Code Analyzer reported no code issues in the new probe; it fell back to
default settings because a local preference file was absent. Unit tests from
the earlier campaign are not reported as newly executed here.

The strongest limitation is that reducing the demonstrated conservatism
does not establish a nonempty robust safe set for every current sensor
enclosure. Even a mathematically correct robust controller can have no
common safe input when admissible target motions demand incompatible
actions. The current measurements establish explicit implementation-induced
contradictions; they do not prove global physical avoidability or the
sufficiency of this proposed redesign. The additional nominal exact-state
failures must also be addressed.

\textit{Research method and disclosure.} This AI-assisted analysis combines
source inspection, saved-snapshot counterfactual evidence, a new nonlinear
terminal probe, explicit algebra and verification against a primary research
paper. The design was checked against the counterarguments that nominal
stability might suffice, that all failures came from target uncertainty, and
that smaller damping must reduce mismatch to zero. Each is contradicted by
the cited local evidence or the stated limit calculation. No independent
experimental validation or implemented robust-controller success is claimed.

\begin{thebibliography}{9}
\bibitem{Kohler} J. K\"ohler, M. A. M\"uller and F. Allg\"ower,
``Robust output feedback model predictive control using online estimation
bounds,'' 2021 preprint, Sections III--IV, especially IV-A and IV-C.
\href{https://arxiv.org/html/2105.03427v1}{arXiv:2105.03427v1}.
\bibitem{Local} Collision Avoidance Research,
\path{controller/OBSERVER_ROBUST_PCBF_THEORY.tex}, especially the causal
error tube, invariant information family and single-CLF uncertainty supply;
\path{report/PERCEPTION_50M_CAMPAIGN_20261003.tex} and its saved experiment
data. The present report identifies the implementation consequences of that
existing derivation, rather than claiming a new general output-feedback theorem.
\end{thebibliography}
\end{document}
